\chapter{The U-Pb system}
\label{sec:U-Pb}

U and Th are found on the extremely heavy end of the Periodic Table of
Elements.  All their isotopes are radioactive and exhibit
$\alpha$-decay and sometimes even spontaneous fission (see Section
\ref{sec:fission-tracks}). $^{232}$Th, $^{235}$U and $^{238}$U each
form the start of long decay series comprising multiple $\alpha$- and
$\beta$ emissions which eventually produce various isotopes of Pb:

\begin{equation}
\begin{array}{rl}
^{238}U \rightarrow & {}^{206}Pb + 8\alpha + 6\beta + 47\mbox{MeV} \\ 
^{235}U \rightarrow & {}^{207}Pb + 7\alpha + 4\beta + 45\mbox{MeV} \\
^{232}Th \rightarrow & {}^{208}Pb + 6\alpha + 4\beta + 40\mbox{MeV} 
\end{array}
\label{eq:UThdecay}
\end{equation}

Each of these three decay series is unique, i.e. no isotope occurs in
more than one series (Figure \ref{fig:U-Th-series}). Furthermore, the
half life of the parent isotope is much longer than any of the
intermediary daughter isotopes, thus fulfilling the requirements for
secular equilibrium (Section \ref{sec:decay-series}). We can therefore
assume that the $^{206}$Pb is directly formed by the $^{238}$U, the
$^{207}$Pb from the $^{235}$U and the $^{208}$Pb from the
$^{232}$Th. Several chronometers are based on the $\alpha$-decay of U
and Th:

\begin{itemize}
\item The U-Th-Pb method (Section \ref{sec:U-Th-Pb})
\item The Pb-Pb method (Section \ref{sec:Pb-Pb})
\item The U-Th-He method (Section \ref{sec:U-Th-He})
\end{itemize}

\section{The U-(Th-)Pb method}
\label{sec:U-Th-Pb}

Natural Pb consists of four isotopes $^{204}$Pb, $^{206}$Pb,
$^{207}$Pb and $^{208}$Pb. The ingrowth equations for
the three radiogenic Pb isotopes are given by:

\begin{equation}
  \begin{array}{rl}
    &{}^{206}Pb^* = {}^{238}U \left(e^{\lambda_{238}t} - 1\right)\\ 
    &{}^{207}Pb^* = {}^{235}U \left(e^{\lambda_{235}t} - 1\right)\\ 
    &{}^{208}Pb^* = {}^{232}Th \left(e^{\lambda_{232}t} - 1\right)
  \end{array}
  \label{eq:Pb*}
\end{equation}

With $\lambda_{238}$ = 1.55125 $\times 10^{-10}$ yr$^{-1}$ (t$_{1/2}$
= 4.468 Gyr), $\lambda_{235}$ = 9.8485 $\times 10^{-10}$ yr$^{-1}$
(t$_{1/2}$ = 703.8 Myr), and $\lambda_{232}$ = 0.495 $\times 10^{-10}$
yr$^{-1}$ (t$_{1/2}$ = 14.05 Gyr). The corresponding age equations
are:

\begin{equation}
  \begin{array}{rl}
    t_{206} & = \frac{1}{\lambda_{238}}
    \ln \left(\frac{{}^{206}Pb^*}{{}^{238}U}+1\right)\\
    t_{207} & = \frac{1}{\lambda_{235}}
    \ln \left(\frac{{}^{207}Pb^*}{{}^{235}U}+1\right)\\
    t_{208} & = \frac{1}{\lambda_{232}}
    \ln \left(\frac{{}^{208}Pb^*}{{}^{232}Th}+1\right)
  \end{array}
  \label{eq:tPb*}
\end{equation}

Some igneous minerals (notably zircon) conveniently incorporate lots
of U and virtually no Pb upon crystallisation. For those minerals, the
non-radiogenic Pb can be safely neglected (at least for relatively
young ages), so that we can assume that $Pb \approx Pb^*$. This
assumption cannot be made for other minerals, young ages, and high
precision geochronology. In those cases, the inherited component (aka
`common Pb') needs to be quantified, which is done by normalising to
non-radiogenic $^{204}$Pb:

\begin{equation}
\begin{array}{c}
  \frac{^{206}Pb}{^{204}Pb} = \left(\frac{^{206}Pb}{^{204}Pb}\right)_\circ +
  \frac{^{238}U}{^{204}Pb} \left(e^{\lambda_{238}t} - 1\right) \\
  \frac{^{207}Pb}{^{204}Pb} = \left(\frac{^{207}Pb}{^{204}Pb}\right)_\circ +
  \frac{^{235}U}{^{204}Pb} \left(e^{\lambda_{235}t} - 1\right)\\
  \frac{^{208}Pb}{^{204}Pb} = \left(\frac{^{208}Pb}{^{204}Pb}\right)_\circ +
  \frac{^{232}Th}{^{204}Pb} \left(e^{\lambda_{232}t} - 1\right)
\end{array}
\label{eq:Pb}
\end{equation}

where $\left(\frac{^{x}Pb}{^{204}Pb}\right)_\circ$ stands for the
common Pb component for isotope x. The corresponding age equations
then become:

\begin{equation}
\begin{array}{c}
  t_{206}=\frac{1}{\lambda_{238}}\ln\left(\frac{\left(\frac{^{206}Pb}{^{204}Pb}\right)-
    \left(\frac{^{206}Pb}{^{204}Pb}\right)_\circ}{\frac{^{238}U}{^{204}Pb}}+1\right)\\
  t_{207}=\frac{1}{\lambda_{235}}\ln\left(\frac{\left(\frac{^{207}Pb}{^{204}Pb}\right)-
    \left(\frac{^{207}Pb}{^{204}Pb}\right)_\circ}{\frac{^{235}U}{^{204}Pb}}+1\right)\\
  t_{208}=\frac{1}{\lambda_{232}}\ln\left(\frac{\left(\frac{^{208}Pb}{^{204}Pb}\right)-
    \left(\frac{^{208}Pb}{^{204}Pb}\right)_\circ}{\frac{^{232}Th}{^{204}Pb}}+1\right)
\end{array}
\label{eq:tPb}
\end{equation}

U-Pb dating grants access to two separate geochronometers
($^{206}$Pb/${}^{238}$U and $^{207}$Pb/${}^{235}$U) based on different
isotopes of the same parent-daughter pair (i.e. U \& Pb).  This
built-in redundancy provides a powerful internal quality check which
makes the method arguably the most robust and reliable dating
technique in the geological toolbox. The initial Pb composition can
either be determined by analysing the Pb composition of a U-poor
mineral (e.g., galena or feldspar) or by applying the isochron method
to samples with different U and Th concentrations. As is the case for
any isotopic system, the system needs to remain `closed' in order to
yield meaningful isotopic ages.  This sometimes is not the case,
resulting in a loss of Pb and/or U.  Such losses cause the
$^{206}$Pb/$^{238}$U- and $^{207}$Pb/$^{235}$U-clocks to yield
different ages. Note that isotopic closure is required for all
intermediary isotopes as well.  Critical isotopes are the highly
volatile $^{226}$Rn (t$_{1/2}$=1.6ka) and $^{222}$Rn
(t$_{1/2}$=3.8d). Initially, the U-Pb method was applied to U-ores,
but nowadays it is predominantly applied to accessory minerals such
zircon and, to a lesser extent, apatite, monazite and allanite.

\section{The Pb-Pb method}
\label{sec:Pb-Pb}

The $^{207}$Pb/$^{206}$Pb method is based on the U-Pb method and is
obtained by dividing the two U-Pb members of Equation \ref{eq:Pb*} (or
\ref{eq:Pb}), and taking into account that the average natural
$^{238}$U/$^{235}$U-ratio is 137.818:

\begin{equation}
\frac{^{207}Pb^*}{^{206}Pb^*} = 
\frac{\left(\frac{^{207}Pb}{^{204}Pb}\right)-\left(\frac{^{207}Pb}{^{204}Pb}\right)_\circ}
{\left(\frac{^{206}Pb}{^{204}Pb}\right)-\left(\frac{^{206}Pb}{^{204}Pb}\right)_\circ}
= \frac{1}{137.818} \frac{e^{\lambda_{235}t}-1}{e^{\lambda_{238}t}-1}
\label{eq:PbPb}
\end{equation}

The left hand side of this equation contains only Pb isotopic ratios.
Note that these are \emph{only} a function of time. Equations
\ref{eq:PbPb} has no direct solution and must be solved
iteratively. The Pb-Pb method has the following advantages over
conventional U-Pb dating:

\begin{itemize}
\item There is no need to measure uranium.
\item The method is insensitive to recent loss of U and even Pb,
  because this would not affect the isotopic ratio of the Pb.
\end{itemize}

In practice, the Pb-Pb method is rarely applied by itself but is
generally combined with the U-Pb technique. The expected
$({}^{207}Pb/{}^{206}Pb)^*$-ratio for recently formed rocks and minerals
can be calculated from Equation \ref{eq:PbPb} by setting
t$\rightarrow$0:

\begin{equation}
  \left(\frac{^{207}Pb}{^{206}Pb}\right)^*_p =
  \frac{\lambda_{235}}{137.818\lambda_{238}} = 0.04607
\label{eq:commonPb}
\end{equation}

This ratio was progressively higher as one goes back further in time.
It was $\approx$ 0.6 during the formation of Earth.

\section{Concordia}
\label{sec:intro2concordia}

It sometimes happens that the U-Th-Pb trio of chronometers does not
yield mutually consistent ages. It is then generally found that
t$_{208}$ $<$ t$_{206}$ $<$ t$_{207}$ $<$ t$_{207/206}$ which, again,
shows that the Pb-Pb clock is least sensitive to open system
behaviour.  From Equation \ref{eq:Pb*}, we find that:

\begin{equation}
\begin{array}{rl}
\frac{^{206}Pb^*}{^{238}U} & = e^{\lambda_{238}t}-1 \mbox{~and}\\
\frac{^{207}Pb^*}{^{235}U} & = e^{\lambda_{235}t}-1
\end{array}
\label{eq:wetherill}
\end{equation}

If we plot those $^{206}$Pb$^*$/$^{238}$U- and
$^{207}$Pb$^*$/$^{235}$U-ratios which yield the same ages (t) against
one another, they form a so-called `concordia' curve. The concordia
diagram is a very useful tool for investigating and interpreting
disruptions of the U-Pb system caused by `episodic lead loss'. This
means that a mineral (of age T$_\circ$, say) has lost a certain
percentage of its radiogenic Pb at a time T$_1$ after its formation
(e.g., during metamorphism), after which the system closes again and
further accumulation of radiogenic Pb proceeds normally until the
present.  On the concordia diagram of multiple aliquots of a sample,
this scenario will manifest itself as a linear array of datapoints
connecting the concordant $^{206}$Pb$^*$/$^{238}$U -
$^{207}$Pb$^*$/$^{235}$U composition expected at T$_\circ$ with that
expected at T$_1$.  With time, the data shift further away from the
origin.  The upper intercept of the linear array (aka \emph{discordia}
line) can be used to estimate the crystallisation age, whereas the
lower intercept yields the age of metamorphism.  The greater the
distance from the expected composition at t, the greater the degree of
Pb loss and the greater the linear extrapolation error on the
crystallisation age (Figure \ref{fig:wetherill}).

\ifpdf
\ifuclnotes
\begin{figure}[!ht]
  \centering
  \def\svgwidth{.7\textwidth}
  \input{wetherill.pdf_tex}
  \caption{`Wetherill' concordia diagram showing concordant (filled
    symbols) and discordant (empty symbols) analyses affected by
    different degrees of Pb (or U) loss \citep[modified
      from][]{allegre2008}.}
  \label{fig:wetherill}
\end{figure}
\else % end of uclnotes
\begin{figure}[!ht]
\noindent\begin{minipage}[t]{.5\textwidth}
\strut\vspace*{-\baselineskip}\newline
\def\svgwidth{\textwidth}
\input{wetherill.pdf_tex}
\end{minipage}
\begin{minipage}[t]{.5\textwidth}
\captionof{figure}{`Wetherill' concordia diagram showing concordant
  (filled symbols) and discordant (empty symbols) analyses affected by
  different degrees of Pb (or U) loss \citep[modified
    from][]{allegre2008}.}
  \label{fig:wetherill}
\end{minipage}
\end{figure}
\fi % end of pdf
\else
\includegraphics[width=10cm]{../figures/wetherill.png}
\captionof{figure}{`Wetherill' concordia diagram showing concordant
  (filled symbols) and discordant (empty symbols) analyses affected by
  different degrees of Pb (or U) loss \citep[modified
    from][]{allegre2008}.}
  \label{fig:wetherill}
\fi

\section{Detrital geochronology}

Zircon (ZrSiO$_4$) is a common U-Th-bearing accessory mineral in
acidic igneous rocks, which form the main proto-sources of the
siliciclastic sediments. Zircon is a very durable mineral that
undergoes minimal chemical alteration or mechanical
abrasion. Therefore, zircon crystals can be considered time capsules
carrying the igneous and metamorphic history of their
proto-sources. The probability distribution of a representative sample
of zircon U-Pb ages from a detrital population can serve as a
characteristic fingerprint that may be used to trace the flow of sand
through sediment routing systems. As a provenance tracer, zircon U-Pb
data are less susceptible to winnowing effects than conventional
petrographic techniques. Using modern microprobe technology (SIMS and
LA-ICP-MS, see Chapter \ref{sec:mass-specs}), it is quite easy to
date, say, a hundred grains of zircon in a matter of just a few
hours. Due to the robustness of zircons as a tracer of sedimentary
provenance, and the relative ease of dating them, the use of detrital
zircon U-Pb geochronology has truly exploded in recent years. A
literature survey using the keywords `detrital', `zircon', and
`provenance' indicates that the proliferation of detrital zircon
studies has followed an exponential trend, with the number of
publications doubling roughly every five years over the past two
decades. At present, nearly a thousand detrital zircon publications
appear each year.\\

An extensive survey of late Archaean sandstones from the Jack Hills in
Australia have revealed a subpopulation of detrital zircons with
Hadean (4.1-4.2 Ga) U-Pb ages. These are the oldest terrestrial
minerals known to science, predating the oldest igneous rocks by 300
million years.  The isotopic composition of oxygen, hafnium and other
elements in the zircon represents a unique window into the earliest
stages of Earth evolution.  They indicate that liquid water was
present on the surface of our planet early on in its history. This
isotopic evidence is corroborated by the geological observation that
the Hadean zircons are preserved in fluvial deposits.
